% Compile with: pdflatex pitch_document.tex
\documentclass[11pt,letterpaper]{article}

\usepackage[T1]{fontenc}
\usepackage{lmodern}
\usepackage[margin=1in]{geometry}
\usepackage{hyperref}
\usepackage{enumitem}
\usepackage{parskip}
\usepackage{xcolor}
\usepackage{fancyhdr}
\usepackage{booktabs}

\hypersetup{
  colorlinks=true,
  linkcolor=blue!70!black,
  urlcolor=blue!60!black,
}

\definecolor{doegreen}{HTML}{005C3B}
\definecolor{doegray}{HTML}{4A4A4A}

\pagestyle{fancy}
\fancyhf{}
\rhead{\textcolor{gray}{\small DOE SBIR/STTR Phase I Proposal}}
\lhead{\textcolor{gray}{\small [Company Name]}}
\rfoot{\textcolor{gray}{\small \thepage}}

\title{%
  \vspace{-1cm}
  \textcolor{doegreen}{\Large\textbf{DOE Office of Science SBIR/STTR Phase I}}\\[0.3cm]
  \textbf{AI-Driven Platform for Accelerated Discovery of\\Advanced Coating Materials for Clean Energy Infrastructure}\\[0.3cm]
  \large [Company Name]\\[0.2cm]
  \normalsize Topic Area: Advanced Materials Discovery Using Machine Learning
}
\author{}
\date{FY 2026}

\begin{document}
\maketitle
\thispagestyle{fancy}

\noindent\textcolor{doegreen}{\rule{\textwidth}{1.5pt}}

\vspace{0.3cm}
\begin{center}
\begin{tabular}{ll}
\textbf{Principal Investigator:} & [PI Name], [Title] \\
\textbf{Organization:} & [Company Name] \\
\textbf{Address:} & [Street Address, City, State ZIP] \\
\textbf{Phone:} & [Phone] \\
\textbf{Email:} & [Email] \\
\textbf{DUNS Number:} & [DUNS] \\
\textbf{Requested Budget:} & \$200,000 (Phase I, 12 months) \\
\textbf{DOE Program:} & Office of Science SBIR/STTR \\
\textbf{Topic:} & Advanced Materials Discovery Using Machine Learning \\
\end{tabular}
\end{center}

\vspace{0.3cm}
\noindent\textcolor{doegreen}{\rule{\textwidth}{1.5pt}}

\vspace{0.5cm}
\noindent\textit{This proposal addresses DOE's mission to accelerate clean energy materials discovery through machine learning. The proposed AI platform will design advanced protective coatings for energy infrastructure---solar panels, wind turbines, hydrogen systems, nuclear facilities, and grid components---reducing materials development timelines from years to weeks while improving durability and sustainability.}

\vspace{0.3cm}
\noindent\textit{Placeholders marked with [brackets] must be replaced with actual company and team information before submission.}

%% =======================================================================
\section{Technical Approach}
\label{sec:technical}
%% =======================================================================

\subsection{Problem Statement: The Materials Bottleneck in Clean Energy Deployment}

The U.S.\ Department of Energy's clean energy goals---net-zero emissions by 2050, 100\% clean electricity by 2035, and a national clean hydrogen economy---depend critically on the durability of deployed infrastructure. Protective coatings are a first line of defense against degradation, yet coating development remains one of the slowest materials R\&D cycles in the energy sector:

\begin{itemize}[noitemsep]
  \item \textbf{Solar photovoltaics:} Encapsulant and anti-reflective coating degradation accounts for 40--60\% of module power losses over 25-year lifetimes. UV-induced yellowing of EVA encapsulants remains a dominant failure mode, costing the U.S.\ solar industry an estimated \$3.5 billion annually in premature replacements.
  \item \textbf{Wind energy:} Leading-edge erosion of turbine blades from rain, hail, and particulates reduces aerodynamic efficiency by 5--25\%, costing \$150,000+ per blade repair. Current polyurethane coatings degrade within 3--5 years offshore.
  \item \textbf{Hydrogen infrastructure:} Hydrogen embrittlement and permeation through pipeline and storage vessel walls demand next-generation barrier coatings. DOE's Hydrogen Shot target of \$1/kg requires dramatically lower infrastructure maintenance costs.
  \item \textbf{Nuclear facilities:} Decontaminable coatings for reactor containment must withstand radiation, chemical exposure, and decontamination cycles. Current qualified coatings (ASTM D5144) have limited lifespans requiring costly shutdowns for recoating.
  \item \textbf{Grid infrastructure:} Corrosion of transmission towers, substations, and underground cables costs U.S.\ utilities \$13.6 billion annually (NACE International). Climate change is accelerating degradation in coastal and high-humidity regions.
  \item \textbf{Geothermal systems:} Downhole equipment coatings must resist temperatures exceeding 250$^\circ$C, highly acidic brines (pH $< 3$), and abrasive particulates---conditions that destroy conventional coatings within months.
\end{itemize}

In every case, the coating material discovery cycle---typically 3--5 years from concept to qualification---is a critical bottleneck. The combinatorial complexity of coating formulation (polymer binders, functional nanoparticle additives, solvent systems, curing conditions) means that traditional Edisonian R\&D explores fewer than 1\% of the feasible design space. DOE national laboratories have recognized this gap: Argonne's Polybot autonomous experimentation platform documented nearly one million fabrication combinations for a single polymer film.

\subsection{Proposed Solution: An Integrated AI Platform for Energy Coating Discovery}

We propose an integrated artificial intelligence platform that unifies three capabilities---generative molecular design, graph neural network (GNN) property prediction, and multi-objective Bayesian optimization (MOBO)---into a closed-loop pipeline purpose-built for advanced coating materials discovery. This platform directly addresses DOE's priorities in ML-accelerated materials science and advanced manufacturing by compressing the coating design cycle from years to weeks.

\subsubsection{Component 1: Energy-Coating-Specific Generative Model}

We will develop a generative molecular model using the Junction Tree Variational Autoencoder (JT-VAE) architecture with a fragment vocabulary specifically curated for energy-relevant coating chemistries:

\begin{itemize}[noitemsep]
  \item \textbf{UV-stabilizing chromophores:} Benzotriazole, hydroxybenzophenone, triazine, and hindered amine light stabilizer (HALS) fragments for solar panel encapsulant and anti-reflective coating design.
  \item \textbf{Corrosion-inhibiting groups:} Phosphonate, silane coupling agent, azole, and rare-earth complexation fragments for anti-corrosion coatings on steel, aluminum, and copper infrastructure.
  \item \textbf{Hydrogen barrier moieties:} Fluoropolymer, polyimide, graphene-functionalized, and clay-nanocomposite fragments for hydrogen permeation barrier coatings.
  \item \textbf{Thermal-stable backbones:} Polysiloxane, polybenzimidazole, polyimide, and ceramic precursor fragments for thermal barrier and geothermal coatings.
  \item \textbf{Cross-linking and self-healing groups:} Epoxide, isocyanate, Diels-Alder adduct, and disulfide exchange fragments enabling autonomous damage repair.
\end{itemize}

Fragment-based generation ensures 100\% chemical validity (versus $\sim$43\% for character-level approaches) and constrains the search to synthetically accessible molecules. The model will be trained on $>$5,000 coating molecules extracted from DOE technical reports, patent databases, and published literature, encoded using SELFIES for guaranteed validity.

\subsubsection{Component 2: Multitask GNN Property Predictors for Energy-Relevant Properties}

We will train multitask graph neural network models predicting properties critical to energy infrastructure coatings:

\begin{itemize}[noitemsep]
  \item \textbf{UV absorption spectrum:} Peak wavelength and breadth of absorption for UV-protective coatings (solar encapsulants, wind blade coatings).
  \item \textbf{Electrochemical impedance:} Impedance modulus at 0.01\,Hz as a proxy for corrosion barrier performance (grid infrastructure, offshore wind, nuclear).
  \item \textbf{Glass transition temperature ($T_g$):} Thermal stability for high-temperature applications (geothermal, turbine, nuclear).
  \item \textbf{Hydrogen permeability:} Diffusion coefficient and solubility for hydrogen barrier coating design.
  \item \textbf{Adhesion energy:} Interfacial strength on steel, aluminum, and glass substrates relevant to energy infrastructure.
  \item \textbf{Weathering resistance:} Accelerated aging performance prediction (critical for 25--30 year solar and wind lifetimes).
\end{itemize}

Transfer learning from the MACE-MP foundation model (pre-trained on $>$150,000 DFT-calculated structures) will enable accurate predictions from fewer than 500 coating-specific training samples per property. Multitask learning shares knowledge across data-rich properties ($T_g$: 1,200+ samples) and data-scarce properties (hydrogen permeability: $<$100 samples).

\subsubsection{Component 3: Constraint-Aware Multi-Objective Bayesian Optimization}

The MOBO module uses BoTorch's qNEHVI (q-Noisy Expected Hypervolume Improvement) acquisition function to navigate the Pareto front across competing objectives. For energy coating applications, we will simultaneously optimize:

\begin{itemize}[noitemsep]
  \item Corrosion resistance (maximize impedance modulus)
  \item UV stability (maximize absorption in target wavelength range)
  \item Thermal stability (maximize $T_g$ or decomposition temperature)
  \item Hydrogen barrier performance (minimize permeability, where applicable)
  \item Mechanical durability (maximize adhesion, minimize brittleness)
\end{itemize}

Subject to constraints including VOC $< 250$\,g/L (EPA regulations), PFAS-free formulation (emerging regulations), and compatibility with existing manufacturing processes (spray coating, dip coating, roll-to-roll deposition). Mixed-variable Gaussian processes with Gower distance kernels handle the combined continuous-categorical design space.

\subsection{Alignment with DOE Office of Science Priorities}

This work directly supports multiple DOE strategic priorities:

\begin{itemize}[noitemsep]
  \item \textbf{Basic Energy Sciences (BES):} Fundamental understanding of structure-property relationships in functional coating materials through data-driven discovery.
  \item \textbf{Energy Efficiency and Renewable Energy (EERE):} Extending the operational lifetime and reducing maintenance costs of solar, wind, and hydrogen infrastructure.
  \item \textbf{Advanced Manufacturing:} Enabling rapid, AI-guided formulation optimization compatible with scalable manufacturing processes.
  \item \textbf{Integrated Computational Materials Engineering (ICME):} Bridging molecular-scale design to macroscopic coating performance through multiscale property prediction.
  \item \textbf{Clean Energy Materials Genome:} Contributing to the Materials Genome Initiative by generating open datasets linking coating molecular structure to energy-relevant performance.
\end{itemize}

%% =======================================================================
\section{Innovation and Impact}
\label{sec:innovation}
%% =======================================================================

\subsection{Scientific Innovation}

The proposed platform introduces three scientific innovations to the field of computational materials discovery:

\textbf{1. First generative AI system purpose-built for energy coating materials.} Existing generative molecular models (GraphVAE, JT-VAE, diffusion models) target drug-like small molecules (20--50 atoms). Coating monomers and oligomers frequently exceed 100 atoms and include functional group combinations absent from pharmaceutical training sets. Our fragment vocabulary, curated from energy coating literature and DOE-funded research, ensures the generative model explores chemically relevant space. No existing commercial or academic tool provides this capability.

\textbf{2. Multiscale property prediction bridging molecular design to coating performance.} Current AI tools for coatings (Citrine Informatics, Uncountable) optimize at the formulation level---adjusting weight percentages of known components---but cannot predict how novel molecular structures will perform. Our GNN models predict molecular-level properties (UV absorption, $T_g$) that directly determine macroscopic coating performance, creating a continuous design loop from atoms to applications.

\textbf{3. Constraint-aware multi-objective optimization for energy-specific requirements.} Energy infrastructure coatings face unique multi-objective challenges: a solar encapsulant must simultaneously maximize UV transmittance in the visible range, block UV below 350\,nm, resist hydrolysis for 25+ years, and maintain adhesion through thermal cycling. Our MOBO framework is the first to handle these energy-specific objective combinations with hard regulatory constraints.

\subsection{Impact on DOE Clean Energy Mission}

\textbf{Accelerated deployment of renewable energy:} By reducing coating development timelines from 3--5 years to 3--6 months, our platform accelerates the qualification of next-generation protective materials for solar, wind, and hydrogen infrastructure. This directly supports DOE's timeline for achieving 100\% clean electricity by 2035.

\textbf{Reduced levelized cost of energy (LCOE):} Coating failures are a major contributor to operations and maintenance costs across energy sectors. Longer-lasting coatings designed by our platform could reduce solar LCOE by 5--8\% (\$0.001--0.002/kWh) and wind LCOE by 3--6\% through extended maintenance intervals.

\textbf{Domestic manufacturing competitiveness:} The platform generates formulations optimized for compatibility with U.S.\ manufacturing processes (spray coating, roll-to-roll, thermal spray), supporting reshoring of advanced coating production and DOE's advanced manufacturing objectives.

\textbf{Sustainability and environmental compliance:} All generated formulations are constrained to comply with EPA VOC limits and emerging PFAS regulations, supporting the transition to environmentally sustainable coating chemistries without sacrificing performance.

\textbf{Open science contributions:} Phase I results, including curated coating property datasets and benchmark model performance, will be published in peer-reviewed journals and deposited in open repositories (e.g., Materials Cloud, NOMAD), consistent with DOE's open science policies.

%% =======================================================================
\section{Phase I Work Plan}
\label{sec:workplan}
%% =======================================================================

Phase I spans 12 months with a total budget of \$200,000. The work is organized into four tasks aligned with DOE materials science priorities.

\subsection{Task 1: Energy Coating Data Curation and Fragment Vocabulary (Months 1--3)}

\textbf{Objective:} Assemble the most comprehensive training dataset of energy-relevant coating molecules and their experimentally measured properties.

\textbf{Activities:}
\begin{itemize}[noitemsep]
  \item Extract molecular structures and property data from DOE technical reports, NIST databases, MatWeb, Reaxys, and patent literature, targeting $>$5,000 unique coating molecules.
  \item Curate a fragment vocabulary of 250+ energy-coating-relevant substructures organized by function (UV absorbers, corrosion inhibitors, hydrogen barriers, thermal stabilizers, cross-linkers).
  \item Compute DFT-level molecular descriptors (HOMO-LUMO gap, dipole moment, polarizability) for molecules lacking electronic structure data, using ORCA at the B3LYP/def2-TZVP level.
  \item Prepare train/validation/test splits stratified by coating application domain (solar, wind, hydrogen, nuclear, grid, geothermal).
\end{itemize}

\textbf{Milestone:} Curated dataset of $\geq$5,000 molecules with $\geq$3 measured properties each; fragment vocabulary validated by domain expert review.

\textbf{Budget:} \$45,000 (personnel, computational resources, database access).

\subsection{Task 2: Generative Model Development and Validation (Months 3--7)}

\textbf{Objective:} Train and validate a generative molecular model that proposes novel coating molecules tailored to energy infrastructure applications.

\textbf{Activities:}
\begin{itemize}[noitemsep]
  \item Implement Junction Tree VAE with energy-coating fragment vocabulary using SELFIES encoding.
  \item Train on curated dataset with augmentation from GEOM-Drugs conformer data for larger molecules ($>$50 atoms).
  \item Evaluate: generate 1,000 novel molecules per energy application domain (6,000 total) and assess chemical validity ($>$99\% target), novelty ($>$80\% not in training set), synthetic accessibility (SA score $< 4.0$ for $>$85\%), and diversity (internal Tanimoto similarity $< 0.4$).
  \item Conduct domain-expert evaluation of top 100 candidates per application for chemical plausibility.
\end{itemize}

\textbf{Milestone:} Generative model producing 6,000+ novel, valid, synthetically accessible coating molecules across six energy application domains.

\textbf{Budget:} \$55,000 (personnel, GPU compute on DOE ASCR resources or commercial cloud).

\subsection{Task 3: Multitask GNN Property Prediction (Months 4--9)}

\textbf{Objective:} Develop GNN property predictors achieving $R^2 > 0.80$ for energy-relevant coating properties.

\textbf{Activities:}
\begin{itemize}[noitemsep]
  \item Implement E(3)-equivariant GNN architecture based on SchNet/DimeNet++ with multitask output heads for six properties (UV absorption, impedance, $T_g$, hydrogen permeability, adhesion, weathering resistance).
  \item Apply transfer learning from MACE-MP foundation model, freezing early message-passing layers and fine-tuning task-specific heads.
  \item Benchmark against random forest baselines on Morgan fingerprints and Gaussian process models.
  \item Quantify prediction uncertainty via deep ensembles (5 models) for downstream Bayesian optimization.
  \item Validate on held-out test sets; target $R^2 > 0.80$ for each property.
\end{itemize}

\textbf{Go/No-Go Decision (Month 8):} If $R^2 < 0.70$ for any property after transfer learning and ensemble approaches, pivot to gradient-boosted tree ensembles with physics-informed molecular descriptors.

\textbf{Milestone:} Validated multitask GNN achieving $R^2 > 0.80$ for $\geq$5 of 6 target properties with calibrated uncertainty estimates.

\textbf{Budget:} \$55,000 (personnel, GPU compute, MACE-MP licensing if applicable).

\subsection{Task 4: Pipeline Integration and Energy Application Demonstrations (Months 8--12)}

\textbf{Objective:} Integrate all components into a closed-loop pipeline and demonstrate AI-guided coating material discovery for two priority DOE energy applications.

\textbf{Activities:}
\begin{itemize}[noitemsep]
  \item Connect generative model, GNN predictors, and MOBO optimizer into a single automated workflow.
  \item \textbf{Demonstration 1---Solar encapsulant coatings:} Optimize for UV blocking (300--380\,nm), visible transmittance ($>$90\%), hydrolysis resistance, and thermal cycling stability. Screen 10,000+ generated candidates and identify Pareto-optimal top 50.
  \item \textbf{Demonstration 2---Anti-corrosion coatings for grid infrastructure:} Optimize for impedance modulus, adhesion on steel, salt spray resistance, and VOC compliance. Screen 10,000+ candidates and identify top 50.
  \item Benchmark pipeline against random screening and single-objective optimization baselines.
  \item Prepare Phase II experimental validation plan with DOE national laboratory and/or industry partners.
\end{itemize}

\textbf{Milestone:} Integrated pipeline demonstrated on two energy applications, producing ranked candidate lists ready for Phase II experimental synthesis and testing.

\textbf{Budget:} \$45,000 (personnel, compute, reporting).

\subsection{Budget Summary}

\begin{center}
\begin{tabular}{lrr}
\toprule
\textbf{Task} & \textbf{Budget} & \textbf{Months} \\
\midrule
Task 1: Data Curation \& Fragment Vocabulary & \$45,000 & 1--3 \\
Task 2: Generative Model Development & \$55,000 & 3--7 \\
Task 3: Multitask GNN Property Prediction & \$55,000 & 4--9 \\
Task 4: Pipeline Integration \& Demonstrations & \$45,000 & 8--12 \\
\midrule
\textbf{Total Phase I} & \textbf{\$200,000} & \textbf{12 months} \\
\bottomrule
\end{tabular}
\end{center}

%% =======================================================================
\section{Commercial Potential}
\label{sec:commercial}
%% =======================================================================

\subsection{Market Size and Growth}

The platform addresses multiple large and growing markets driven by the clean energy transition:

\begin{itemize}[noitemsep]
  \item \textbf{Global protective coatings market:} \$17 billion (2025), projected \$32.4 billion by 2034 (7.3\% CAGR, Precedence Research).
  \item \textbf{Solar encapsulant materials:} \$3.2 billion (2025), growing at 12\% CAGR with global solar deployment.
  \item \textbf{Wind turbine coatings:} \$1.8 billion (2025), driven by offshore wind expansion.
  \item \textbf{Hydrogen infrastructure materials:} Emerging market expected to reach \$2.5 billion by 2030 (DOE Hydrogen Shot initiative).
  \item \textbf{Anti-corrosion coatings for energy infrastructure:} \$8.2 billion (2025), with utilities and grid operators as primary customers.
  \item \textbf{Nuclear-qualified coatings:} \$600 million niche market with extremely high barriers to entry.
\end{itemize}

Our serviceable addressable market (SAM) is the high-performance energy infrastructure coatings segment, estimated at \$4.5 billion, where formulation complexity and performance requirements justify AI-driven optimization.

\subsection{Commercialization Strategy}

\textbf{Phase I $\rightarrow$ Phase II Transition:} Phase I delivers a validated software prototype. Phase II (\$1.5M, 24 months) will fund experimental synthesis and testing of top candidates at DOE national laboratory partner facilities, targeting 3--5 experimentally validated coating formulations per energy application domain.

\textbf{Revenue Model:}
\begin{itemize}[noitemsep]
  \item \textbf{SaaS Platform:} Annual licensing at \$150K--300K per seat for coating manufacturers and energy companies. Premium tiers include custom model training on proprietary formulation data.
  \item \textbf{Discovery Partnerships:} Co-development agreements with coating manufacturers (PPG, Hempel, Sherwin-Williams) and energy technology companies (First Solar, Vestas, NextEra Energy), with milestone payments and royalties on commercialized formulations.
  \item \textbf{Government Contracts:} Follow-on DOE, DOD, and NASA contracts for mission-specific coating optimization (e.g., extreme environment coatings for fusion reactors, space-rated coatings).
\end{itemize}

\textbf{Intellectual Property:} Phase I will generate IP in three areas: (1) energy-coating fragment vocabularies and training datasets, (2) trained GNN models for energy-relevant property prediction, and (3) the integrated pipeline architecture. Patent applications will be filed for novel coating molecules identified through the platform.

\subsection{Customer Validation}

[Company Name] has initiated discussions with potential Phase II partners including:
\begin{itemize}[noitemsep]
  \item [Coating Manufacturer] --- interest in AI-accelerated development of PFAS-free anti-corrosion coatings.
  \item [Energy Company/Utility] --- interest in next-generation solar encapsulant materials with extended UV stability.
  \item [DOE National Laboratory] --- potential collaboration on experimental validation using [Lab's] coating characterization facilities.
\end{itemize}

\textbf{Letters of interest from industry and national laboratory partners will be included in the full proposal.}

%% =======================================================================
\section{Related Experience and Qualifications}
\label{sec:team}
%% =======================================================================

\subsection{Principal Investigator}

\textbf{[PI Name]}, [Title], [Company Name]. [PI Name] brings [X] years of experience in computational materials science, machine learning, and molecular design. [Relevant credentials: PhD in chemistry, materials science, chemical engineering, or computer science with materials focus. Publications in peer-reviewed journals on GNN property prediction, generative molecular design, or computational coatings science. Prior experience with DOE-funded research or national laboratory collaborations.]

Relevant expertise includes:
\begin{itemize}[noitemsep]
  \item Graph neural network development for materials property prediction
  \item Generative models for molecular design (VAE, diffusion models, reinforcement learning)
  \item Multi-objective optimization for materials discovery
  \item [Specific coating or polymer materials experience]
\end{itemize}

\subsection{Co-Investigators and Key Personnel}

\textbf{[Co-PI Name]}, [Title]. Brings [X] years of expertise in protective coating chemistry, polymer formulation, and experimental characterization. [Relevant credentials: PhD or MS in polymer science, coatings technology, or corrosion engineering. Industrial R\&D experience at coating manufacturers or energy companies. Expertise in EIS, UV-Vis spectroscopy, salt spray testing, accelerated weathering.]

\textbf{[ML Engineer Name]}, [Title]. [X] years of experience in deep learning infrastructure, model deployment, and scientific software development. [Relevant credentials: experience with PyTorch, PyTorch Geometric, BoTorch; prior ML platform development.]

\subsection{Organizational Capabilities}

[Company Name], founded [Year] in [City, State], is dedicated to replacing Edisonian coating R\&D with AI-driven materials discovery. The company maintains:

\begin{itemize}[noitemsep]
  \item Computational infrastructure: [GPU cluster specifications or cloud computing agreements for model training].
  \item Software stack: Production-grade Python/PyTorch/BoTorch pipeline with version control, CI/CD, and automated testing.
  \item Research partnerships: Collaboration agreements with [University] providing access to [characterization facilities] and with [DOE National Laboratory] for [specific capabilities].
  \item Advisory board: [Names and affiliations of advisors in computational materials science, coating industry, and energy sector].
\end{itemize}

\subsection{Prior Funding and Relevant Experience}

[If applicable: Prior SBIR/STTR awards from NSF, DOE, DOD, or other agencies. Description of related prior work, publications, or preliminary results.]

[If new company: Description of founders' relevant prior work at universities, national laboratories, or industry, including publications and patents related to AI-driven materials discovery.]

%% =======================================================================
\section{References}
\label{sec:references}
%% =======================================================================

\begin{enumerate}[noitemsep]

\item Armelin, E., et al.\ ``Corrosion protection with polyaniline and polypyrrole as anticorrosion coatings.'' \textit{Corrosion Science} 50, 721--728 (2008).

\item Axelrod, S. \& Gomez-Bombarelli, R.\ ``GEOM: Energy-annotated molecular conformations.'' \textit{Scientific Data} 9, 185 (2022).

\item Baxter, D.\ ``DOE hydrogen infrastructure materials challenges.'' DOE Hydrogen Program Review (2024).

\item Degen, J., et al.\ ``On the art of compiling and using `drug-like' chemical fragment spaces.'' \textit{ChemMedChem} 3, 1503--1507 (2008).

\item Deshwal, A., et al.\ ``Bayesian optimization over high-dimensional combinatorial spaces via dictionary-based embeddings.'' \textit{NeurIPS} (2023).

\item Dunn, A., et al.\ ``Benchmarking materials property prediction methods.'' \textit{npj Computational Materials} 6, 138 (2020).

\item Hachmann, J., et al.\ ``The Harvard Clean Energy Project: large-scale computational screening and design of organic photovoltaics.'' \textit{Energy \& Environmental Science} 4, 4849--4861 (2011).

\item Hoogeboom, E., et al.\ ``Equivariant diffusion for molecule generation in 3D.'' \textit{ICML} (2022).

\item Jain, A., et al.\ ``Commentary: The Materials Project.'' \textit{APL Materials} 1, 011002 (2013).

\item Jin, W., et al.\ ``Junction tree variational autoencoder for molecular graph generation.'' \textit{ICML} (2018).

\item Koch, G., et al.\ ``International measures of prevention, application, and economics of corrosion technologies study.'' NACE International (2016).

\item Krenn, M., et al.\ ``Self-referencing embedded strings (SELFIES): A 100\% robust molecular string representation.'' \textit{Machine Learning: Science and Technology} 1, 045024 (2020).

\item Musil, F., et al.\ ``Physics-inspired structural representations for molecules and materials.'' \textit{Chemical Reviews} 121, 9759--9815 (2021).

\item National Renewable Energy Laboratory.\ ``Best practices for operation and maintenance of photovoltaic and energy storage systems.'' NREL Technical Report TP-7A40-73822 (2018).

\item Paszke, A., et al.\ ``PyTorch: An imperative style, high-performance deep learning library.'' \textit{NeurIPS} (2019).

\item Pilania, G., et al.\ ``Machine learning bandgaps of double perovskites.'' \textit{Scientific Reports} 6, 19375 (2016).

\item Ramprasad, R., et al.\ ``Machine learning in materials informatics: recent applications and prospects.'' \textit{npj Computational Materials} 3, 54 (2017).

\item Szymanski, N.J., et al.\ ``An autonomous laboratory for the accelerated synthesis of novel materials.'' \textit{Nature} 624, 86--91 (2023).

\item U.S.\ Department of Energy.\ ``Hydrogen Shot: An Introduction.'' DOE Hydrogen and Fuel Cell Technologies Office (2023).

\item Ward, L., et al.\ ``Matminer: An open source toolkit for materials data mining.'' \textit{Computational Materials Science} 152, 60--69 (2018).

\item Xie, T. \& Grossman, J.C.\ ``Crystal graph convolutional neural networks for an accurate and interpretable prediction of material properties.'' \textit{Physical Review Letters} 120, 145301 (2018).

\end{enumerate}

\vspace{1cm}
\noindent\textcolor{doegreen}{\rule{\textwidth}{1.5pt}}

\begin{center}
\small\textcolor{doegray}{
DOE Office of Science SBIR/STTR FY 2026\\
Topic: Advanced Materials Discovery Using Machine Learning\\
[Company Name] --- Confidential
}
\end{center}

\end{document}
